\subsection{Phân rã tensor trong học sâu}

Phần này trình bày cách các cấu trúc tensor xuất hiện một cách tự nhiên trong
học sâu hiện đại và cách phân rã tensor, hay low-rank factorization, được dùng
để giảm bộ nhớ, giảm chi phí tính toán và thiết kế kiến trúc mới. Trong tutorial
\textit{Foundations of Tensor Computations for AI}, phần deep learning được đặt
trong bối cảnh ``AI scaling wall'': mô hình ngày càng lớn, chi phí huấn luyện và
suy diễn ngày càng cao, trong khi khả năng triển khai trên thiết bị nhỏ và tài
nguyên hạn chế ngày càng quan trọng \cite{chrysos2025tutorial}.

\subsubsection{Tensor xuất hiện tự nhiên trong học sâu}

Trong các thư viện như PyTorch, tensor thường được xem như một kiểu dữ liệu để
lưu trữ mảng nhiều chiều. Tuy nhiên, góc nhìn của tutorial mạnh hơn: hình dạng
nhiều chiều của dữ liệu và tham số phản ánh cấu trúc của bài toán, và cấu trúc
đó có thể được khai thác bằng các phép toán multilinear.

Một batch ảnh có dạng
\begin{equation}
    \mathcal{X} \in \mathbb{R}^{B \times C \times H \times W},
\end{equation}
trong đó $B$ là batch size, $C$ là số kênh, $H,W$ là chiều cao và chiều rộng.
Trong mạng tích chập, trọng số của một lớp convolution là tensor bậc 4:
\begin{equation}
    \mathcal{W} \in \mathbb{R}^{C_{\mathrm{out}} \times C_{\mathrm{in}} \times K_H \times K_W}.
\end{equation}
Mỗi mode của tensor này có ý nghĩa riêng: kênh đầu ra, kênh đầu vào, và cấu trúc
không gian của kernel.

Trong Transformer, multi-head attention cũng là một phép tính tensor. Với
$Q,K,V \in \mathbb{R}^{B \times H \times L \times d}$, trong đó $H$ là số head,
$L$ là độ dài chuỗi và $d$ là head dimension, attention score có dạng
\begin{equation}
    A_{bhij} = \sum_k Q_{bhik}K_{bhjk},
    \qquad
    A \in \mathbb{R}^{B \times H \times L \times L}.
\end{equation}
Output của attention được tính bằng
\begin{equation}
    O_{bhi\delta} = \sum_j A_{bhij}V_{bhj\delta}.
\end{equation}
Hai biểu thức này cho thấy attention không chỉ là nhân ma trận thông thường:
đó là một chuỗi tensor contractions có cấu trúc theo batch, head, token và
feature dimension \cite{vaswani2017attention}.

\subsubsection{Polynomial networks và tương tác bậc cao}

Tutorial tiếp tục từ nhận xét rằng neural network cần học các tương tác phức tạp
giữa đặc trưng đầu vào. Một cách mô hình hóa trực tiếp là viết hàm cần học dưới
dạng khai triển đa thức. Với hàm $G:\mathbb{R}^{d}\rightarrow\mathbb{R}^{o}$,
một mô hình bậc $N$ có thể viết dưới dạng:
\begin{equation}
    G(x) =
    \beta
    + \mathcal{W}^{[1]} \times_2 x
    + \mathcal{W}^{[2]} \times_2 x \times_3 x
    + \cdots
    + \mathcal{W}^{[N]} \times_2 x \cdots \times_{N+1} x.
\end{equation}
Trong đó, $\mathcal{W}^{[n]}$ là tensor trọng số nắm bắt tương tác bậc $n$ giữa
các bản sao của input. Cách viết này giải thích vì sao tensor là đối tượng tự
nhiên khi mô hình cần biểu diễn high-order interactions.

Vấn đề là số tham số của $\mathcal{W}^{[n]}$ tăng theo lũy thừa của $d$.
Pi-nets giải quyết bằng cách tham số hóa các tensor trọng số dưới dạng phân rã
low-rank, đặc biệt gắn với CP decomposition và tích Hadamard
\cite{chrysos2020pi,chrysos2022augmenting}. Thay vì lưu toàn bộ tensor đầy đủ,
mô hình học các factor matrices có kích thước nhỏ hơn. Đây là ví dụ đầu tiên cho
thông điệp chính của tutorial: nếu weight tensor có cấu trúc, ta có thể học biểu
diễn gọn hơn mà vẫn giữ được khả năng mô hình hóa.

\subsubsection{Mixture of Experts dưới góc nhìn tensor}

Mixture of Experts (MoE) là một ví dụ khác về việc trọng số của mô hình có thể
được nhìn như tensor. Một linear layer thông thường tính
\begin{equation}
    y = W^\top x.
\end{equation}
Trong MoE, ta có $N$ expert với các ma trận $W_1,\ldots,W_N$ và hệ số gating
$a_n(x)$ phụ thuộc đầu vào:
\begin{equation}
    y = \sum_{n=1}^{N} a_n(x) W_n^\top x.
\end{equation}
Nếu xếp các expert weights lại với nhau, ta thu được tensor
\begin{equation}
    \mathcal{W} \in \mathbb{R}^{N \times I \times O}.
\end{equation}
Góc nhìn này cho phép áp dụng tensor factorization vào MoE để chia sẻ tham số
giữa các expert, giảm kích thước mô hình và vẫn cho phép expert specialization.
Multilinear Mixture of Experts khai thác trực tiếp cấu trúc này để mở rộng MoE
theo hướng factorized và hiệu quả hơn \cite{oldfield2024multilinear}.

\subsubsection{LoRA và Low-Rank Tensor Adaptation}

Trong fine-tuning truyền thống, mỗi ma trận trọng số $W_0 \in
\mathbb{R}^{k \times d}$ có thể nhận một update đầy đủ $\Delta W$ cùng kích
thước. Cách này linh hoạt nhưng rất tốn tham số, đặc biệt với language models
lớn. LoRA đề xuất đóng băng $W_0$ và chỉ học một update low-rank
\cite{hu2022lora}:
\begin{equation}
    h = W_0x + \frac{\alpha}{r}BAx,
    \qquad
    A \in \mathbb{R}^{r \times d},\quad
    B \in \mathbb{R}^{k \times r}.
\end{equation}
Số tham số trainable giảm từ $kd$ xuống
\begin{equation}
    r(d+k) \ll kd
\end{equation}
khi rank $r$ nhỏ hơn nhiều so với $d,k$. Đây là lý do LoRA phù hợp với demo của
đồ án: ta có thể đo trực tiếp trade-off giữa validation accuracy và số tham số
trainable khi thay đổi rank $r$.

LoRTA mở rộng tư duy này từ ma trận sang tensor adaptation: các update của nhiều
layer và nhiều attention heads được xem như một đối tượng bậc cao, sau đó
factor hóa và chia sẻ factor giữa các mode \cite{hounie2024lorta}. Nói cách
khác, LoRA là trường hợp matrix low-rank adaptation, còn LoRTA đưa ý tưởng đó
gần hơn với tensor factorization đúng nghĩa.

\subsubsection{TensorGRAD và attention hiệu quả}

Tutorial cũng đưa ra hai hướng deep learning khác. TensorGRAD xem gradient trong
huấn luyện neural operators như một tensor bậc cao và dùng Tucker decomposition
để giảm chi phí bộ nhớ \cite{loeschcke2025tensorgrad}. Thay vì lưu và cập nhật
gradient đầy đủ, phương pháp làm việc với core tensor và sparse factors, phù hợp
với các setting mà gradient quá lớn.

Với attention, self-attention chuẩn có chi phí theo độ dài chuỗi gần $O(L^2)$.
Poly-SA thay thế một phần cấu trúc softmax attention bằng các phép tính dựa trên
tích Hadamard:
\begin{equation}
    f_{\mathrm{Poly\mbox{-}SA}}(X)
    =
    \Psi\big((XW_1) * (XW_2)\big) * XW_3,
\end{equation}
hướng tới chi phí gần tuyến tính theo số token \cite{babiloni2023polysa}. Gần
đây, Tensor Product Attention tiếp tục khai thác tensor product để thiết kế lại
attention blocks \cite{zhang2025tensorproductattention}. Các ví dụ này cho thấy
tensor không chỉ dùng để nén tham số, mà còn có thể dẫn tới kiến trúc mới.

\subsubsection{Liên hệ với demo}

Trong phần thực nghiệm, nhóm hiện thực demo \textit{LoRA rank sweep} trên bài
toán phân loại cảm xúc. Demo so sánh frozen encoder, full fine-tuning và LoRA
với các rank khác nhau. Các metric chính gồm validation accuracy, validation
loss, số tham số trainable và thời gian huấn luyện. Kết quả kỳ vọng không phải
là LoRA luôn vượt full fine-tuning, mà là LoRA có thể đạt trade-off tốt hơn giữa
hiệu năng và số tham số cần cập nhật. Điều này bám sát thông điệp của tutorial:
khai thác cấu trúc low-rank/multilinear giúp mô hình học sâu trở nên gọn và dễ
triển khai hơn.
